\section{Related Work}

\subsection{Static and Symbolic Analysis}
The foundational layer of smart contract security relies on static analysis and symbolic execution. Tools such as \textbf{Oyente}, \textbf{Mythril}, and \textbf{Slither} analyze contract bytecode or source code to detect known vulnerability patterns like Reentrancy, Integer Overflows, and Unchecked External Calls. \textbf{Securify} introduces formal verification to prove the violation of safety properties. 
While effective for identifying coding errors during the development phase, these tools suffer from significant limitations in the DeFi era:
\begin{enumerate}
    \item \textbf{Logic Blindness:} They cannot detect business logic flaws or "Rug Pull" mechanisms (e.g., a hidden minting function) that are syntactically correct but semantically malicious.
    \item \textbf{Lack of Runtime Context:} Static tools cannot analyze transaction sequences or external state dependencies (e.g., Flash Loan attacks), rendering them ineffective against dynamic exploits.
\end{enumerate}
TrustManager differs by focusing on \textit{runtime behavioral analysis}, capturing malicious intent through transaction execution traces rather than code structure.

\subsection{Machine Learning in Blockchain Security}
Recent literature has explored Machine Learning (ML) for detecting anomalies in blockchain networks. 
\textbf{Ponzi Scheme Detection:} Early works utilized feature engineering (e.g., Gini coefficient of transaction values) with XGBoost or SVM to identify Ponzi schemes on Ethereum.
\textbf{Phishing & Fraud Detection:} More advanced approaches employ Graph Neural Networks (GNNs) or Random Walks (e.g., \textit{Trans2Vec}) to model the transaction graph and propagate guilt from known malicious nodes.
\textbf{Limitations of Prior ML Works:} Most existing studies operate in a \textit{post-mortem} or \textit{offline} manner. They analyze historical data to classify addresses long after the attack has occurred. Furthermore, deep learning models (GNN/LSTM) often lack explainability, making them unsuitable for high-stakes financial decisions where users demand to know \textit{why} a transaction was blocked. Our work prioritizes \textit{explainable AI} (Random Forest with feature importance) and \textit{real-time} prevention.

\subsection{Graph Learning and Explainability}
Graph Neural Networks have become a dominant approach for modeling relational structures on blockchains due to their ability to propagate signals along edges and capture higher-order interactions. Comprehensive surveys on GNNs summarize architectures and training challenges in large, dynamic graphs. While GNNs achieve high detection accuracy, explainability remains an open problem for high-stakes finance; therefore, models with transparent importance attribution (e.g., Random Forests with permutation importance) remain attractive for user-facing systems. Classical graph-theoretic measures such as PageRank and degree centrality are effective lightweight alternatives that preserve interpretability while capturing global and local influence.

\subsection{Real-time Monitoring & Security Oracles}
Operational security tools like \textbf{Forta}, \textbf{OpenZeppelin Defender}, and \textbf{GoPlus Security} provide real-time monitoring capabilities.
\begin{itemize}
    \item \textbf{Forta} utilizes a decentralized network of bots to scan transactions against predefined heuristics (e.g., "alert if gas usage $>$ 10M").
    \item \textbf{GoPlus} offers an API for token security detection.
\end{itemize}
However, these systems largely rely on \textbf{rigid rule-based heuristics}, which suffer from high false-positive rates (alert fatigue) and are easily bypassed by stealthy attacks (e.g., splitting a large hack into multiple small transactions). Moreover, they typically function as "Alerting Systems" rather than "Prevention Systems"—they notify the user, but often too late to stop the transaction.
TrustManager bridges this gap by introducing a \textbf{Hybrid AI-Oracle} architecture. Unlike passive monitors, our system acts as an active gatekeeper, using consensus-validated ML scores to programmatically block malicious transactions on-chain before finalization.

\subsection{Decentralized Oracles and Rollups}
Decentralized Oracle Networks (DONs) such as Chainlink use threshold signatures and staking to provide tamper-resistant data feeds. Optimistic systems including Optimistic Rollups and optimistic oracles (e.g., UMA) achieve scalability by assuming honest execution and enabling fraud proofs within a challenge window. Our design adopts this paradigm: threshold consensus offers fast-path finality for routine transactions, while a challenge period provides a cryptographic and economic backstop against majority collusion.

\subsection{Gap Analysis}
Table \ref{tab:comparison} summarizes the distinction between TrustManager and existing solutions. We are the first to combine off-chain ML flexibility with on-chain enforceability in a cross-chain context.

\begin{table}[htbp]
\centering
\caption{Comparison with State-of-the-Art Solutions}
\label{tab:comparison}
\begin{tabular}{lcccc}
\toprule
\textbf{Solution} & \textbf{Analysis Type} & \textbf{Detection Logic} & \textbf{Real-time Blocking} & \textbf{Cross-chain} \\
\midrule
Mythril / Slither & Static & Symbolic / Patterns & No (Audit only) & No \\
Forta & Dynamic & Rule-based Bots & No (Alert only) & Yes \\
GNN-based Methods & Graph Analysis & Deep Learning & No (Offline) & No \\
\textbf{TrustManager (Ours)} & \textbf{Hybrid} & \textbf{ML (Random Forest)} & \textbf{Yes (via Oracle)} & \textbf{Yes} \\
\bottomrule
\end{tabular}
\end{table}
